\section{Every measurement and its bound, in one place}\label{sec:G}

So that nothing in this report can be mistaken for measured. What was a fourteen-row table is seven
rows: the six requirement-level experiments that filled it have run and are reported in §\ref{sec:11.2}, and
a seventh row records an experiment this report did not originally plan. The occupancy sensitivity
the first asked for was already in T-24's own result file, R-5's rebuilt case and T-23's dashboard
frames were captured from a live demo-profile run, T-27's kills were repeated against
\texttt{docker compose}, T-15's confusion matrix has been measured over 218 runs, T-29's optional
contender has been measured and dropped, and T-35's reordering of the diagnosis has been measured
and not adopted, so \textbf{every row in this table is now filled}, and each keeps the bound on
what its evidence actually is. That is not the same as every
requirement passing: R-1 and R-3 still fail on the arms named in §\ref{sec:11.2}, R-2 and R-7 fail on the
harder of their two occupancy regimes, and T-15 reports that one of four injected mechanisms is
classified rather than all of them. The table being full means nothing in this report is waiting on a measurement, not
that the measurements came out well.

One result is worth stating separately, because the obvious way to test it measures the wrong
thing. All \factComposeServices{} services declare \texttt{restart: unless-stopped}. After \texttt{docker kill} both
killed containers sat at exit 137 with a restart count of \textbf{zero}, which looks like a broken
policy and is not: a kill issued through the daemon is recorded as a manual stop and deliberately
suppresses restarts. Nor does \texttt{SIGKILL} work from inside, because PID 1 in these images is
the JVM itself with no init wrapper, and the kernel ignores \texttt{SIGKILL} sent to PID 1 from
within its own namespace. The signal that does work is \texttt{SIGTERM} delivered inside the
container: Quarkus handles it, the JVM shuts down of its own accord, and Docker sees a process that
exited rather than one that was stopped. Measured that way the policy fires: exit code 0,
\textbf{restart count 1}, and the container back in \texttt{running} within 3 s. \textbf{The
distinction is the finding}: a chaos test that only runs \texttt{docker kill} cannot tell a working
restart policy from a missing one, and would have reported this configuration as broken.

\begin{longtable}[]{@{}
  >{\raggedright\arraybackslash}p{(\columnwidth - 6\tabcolsep) * \real{0.0878}}
  >{\raggedright\arraybackslash}p{(\columnwidth - 6\tabcolsep) * \real{0.0878}}
  >{\raggedright\arraybackslash}p{(\columnwidth - 6\tabcolsep) * \real{0.4487}}
  >{\raggedright\arraybackslash}p{(\columnwidth - 6\tabcolsep) * \real{0.3707}}@{}}
\toprule
\textbf{Req.} & \textbf{Run} & \textbf{The number that must replace the marker} &
\textbf{Where it goes in this report} \\
\midrule
\endhead
R-1 & T-24 & \textbf{Filled, with a bound on what it is.} Detection delay against the
actual-to-design occupancy ratio is now read off the closed-loop arm's own per-run head count:
0 of 4 confirmed at 0.625 of design, then 122 and 257 min at 0.75, 77 min median at 0.875 and
31 min median at 1.00 (§\ref{sec:11.2}). It is a stratification of 20 drawn weekdays rather than a sweep of
a controlled occupancy, over one room and one $h$, so a true sensitivity experiment is still a
different run & §\ref{sec:7.2}, §\ref{sec:11.2} \\
R-4 & T-27 & \textbf{Filled, except for the restart policy.} The kills were carried out against
\texttt{docker compose} on \texttt{daa2d5c}. Stopping \texttt{co2-sensor} put five of six rooms in
\texttt{BLIND} within 12 s; the sixth was latched and stayed \texttt{COMPENSATING}, which is the
design. No new diagnosis was raised while blind, and no healthy room was commanded. The only
commands issued were the latched room's held position and the fan. Recovery ran
\texttt{BLIND}~→~\texttt{WARMUP}~→~\texttt{HEALTHY} across all five rooms in three cycles.
Killing \texttt{fault-monitor} left ventilation running; killing \texttt{supervisor} mid-compensation
left the fan at 158.3~W, raised no new strategy and kept A115 latched. The restart policy is
measured too, once the kill is delivered in a way Docker counts: see §\ref{sec:11.6} & §\ref{sec:3} status, §\ref{sec:11.2} \\
R-5 & T-20 & \textbf{Filled, with a bound on what it is.} A diagnosis was rebuilt from the journal
live, in the recorded demo run: \texttt{GET /api/cases/aecd9137-\allowbreak 6c3d-\allowbreak 42a8-\allowbreak ac11-\allowbreak 438ef3d9d2d6}
returned HTTP 200, 77\,248 bytes and \textbf{445 events} (every reading, residual and CUSUM value
in sequence) against a real TimescaleDB rather than a test double
(\texttt{evaluation/results/T-23-demo-run.txt}, §\ref{sec:5}). The \texttt{caseId} was the one that run produced rather than a
randomly chosen one. \textbf{R-5 now has an integration test as well}: eight tests execute
\texttt{replayCase}'s real SQL against a real PostgreSQL, including the one-hour evidence lookback,
and they bound reconstruction fidelity at half a microsecond because a \texttt{timestamptz} round
trip rounds rather than truncates & §\ref{sec:3} status \\
R-1, R-3 & T-15 & \textbf{Filled, and it is the least flattering result in this report.} 218 runs
over the four injectable faults plus a no-fault arm. \texttt{STUCK\_DAMPER} is named correctly;
\texttt{BLOCKED\_FILTER} confirms as \texttt{UNCLASSIFIED} 11 of 11, as predicted;
\texttt{SENSOR\_DRIFT} is \textbf{never emitted at all} (0 of 76 runs from a drift ramp over
5--160 ppm/h \emph{and} 0 of 76 from a step offset over 50--700 ppm), so \textbf{one of four
mechanisms is classified}. The no-fault arm is silent on 10 of 11 seeds, and the one false positive
arrives as a \emph{named} damper in a \emph{named} room rather than as \texttt{UNCLASSIFIED}, which
is a worse failure than a vague one (§\ref{sec:7.3}) & §\ref{sec:7.3} \\
R-1, R-3 & T-29 & \textbf{Filled: measured, and dropped.} A one-class detector fitted on a
fault-free commissioning window only, scored on a held-out 7-day week over 20 seeds and 30
operating points. At the incumbent's 0.75 detection rate it costs about \textbf{20× the false
alarms} (11.25 against 0.55 per week) and \textbf{7× the delay} (518 against 74 min), and no point
on the grid beats the monitor. Its misses are
\emph{exactly} the monitor's (the same five seeds), which is the finding: what limits detection
here is the evidence the residual carries, not the decision rule, and a learned rule cannot add
evidence (§\ref{sec:11.2}) & §\ref{sec:11.2}, §\ref{sec:13} \\
— & T-23 & \textbf{Filled, and re-captured after the room-layer fix.} Seven frames from a live
demo-profile run: the dashboard with its detector deliberately not firing at simulated 10:01, the
same table with A115 \texttt{CONFIRMED} at 10:23, BuildSim's floor plan beside the supervisor's
work order at that instant, and \textbf{all four layers at one instant}, the last of which could
not be taken at all before the collection defect was fixed, because three of the four never
reached BuildSim. No requirement states what the view must show, so this is evidence the operator
path works rather than evidence against a criterion & §\ref{sec:12} \\
R-1 & T-35 & \textbf{Filled by an experiment this report did not originally plan, and the answer is
no.} The stuck-damper leaf accepts \texttt{NOT\_JUDGED} because a probe can only be requested once a
case is latched, so the obvious remedy is to withhold the diagnosis until a probe has landed. Built,
and measured over 176 paired runs on identical seeds: detection moves in \textbf{0 of 88 pairs}, the
fault is named on a genuinely judged step in 28 of the 39 runs that defer against 0 before, and
the settled classification is \textbf{identical in all 88 pairs}, because the judged step returns
\texttt{MISSING} for every condition, including nine rooms whose dampers are in perfect health and
one run with nothing injected. \textbf{Not adopted} (§\ref{sec:7.3}, §\ref{sec:13.6}) & §\ref{sec:7.3}, §\ref{sec:10}, §\ref{sec:13.6} \\
\bottomrule
\end{longtable}
